% BHU PYQ -- Auto-generated & error-corrected
\documentclass[11pt,a4paper]{article}

\usepackage[a4paper,top=2.5cm,bottom=2.5cm,left=2.8cm,right=2.8cm,headheight=14pt]{geometry}
\usepackage{amsmath,amssymb}
\usepackage[T1]{fontenc}
\usepackage[utf8]{inputenc}
\usepackage{lmodern}
\usepackage{microtype}
\usepackage{fancyhdr}
\usepackage{enumitem}
\usepackage{hyperref}
\usepackage{lastpage}
\usepackage{parskip}

\hypersetup{colorlinks=true,urlcolor=black,linkcolor=black,
  pdftitle={STA/STB-501: Applied Statistics -- BHU PYQ},pdfauthor={Banaras Hindu University}}

\pagestyle{fancy}\fancyhf{}
\renewcommand{\headrulewidth}{0.4pt}\renewcommand{\footrulewidth}{0.4pt}
\fancyhead[L]{\small   \textit{Banaras Hindu University}}
\fancyhead[R]{\small   \textit{STA/STB-501: Applied Statistics}}
\fancyfoot[C]{\small Page  \thepage\ of \pageref{LastPage}}


\newlist{parts}{enumerate}{2}
\setlist[parts,1]{label=(\alph*),leftmargin=*,itemsep=5pt,topsep=3pt}
\setlist[parts,2]{label=(\roman*),leftmargin=*,itemsep=3pt,topsep=2pt}

\newcommand{\pts}[1]{\hfill   \textit{[#1]}}


\begin{document}


\begin{center}
  {\large\bfseries BANARAS HINDU UNIVERSITY}\\[2pt]
  {\normalsize B.A./B.Sc. (Hons.) Semester V Examination, 2022-23}\\[6pt]
  \rule{0.8\linewidth}{0.5pt}\\[6pt]
  {\large\bfseries Paper No. STA/STB-501: Applied Statistics}\\[4pt]
  {\normalsize Time: 3 Hours\hfill Maximum Marks: 70}
\end{center}
\rule{\linewidth}{0.4pt}
\smallskip



\noindent   \textit{Instructions: Answer any   \textbf{five} questions. All questions carry equal marks. Symbols have their usual meanings.}
\bigskip

\medskip


\noindent   \textbf{Question 1.}

\begin{parts}
  \item What are the methods of obtaining Vital Statistics? Also give the uses of Vital Statistics. \pts{7}
  \item Define Gross Reproduction Rate (GRR) and Net Reproduction Rate (NRR). Explain how far they can be considered as indices of population growth. Interpret the following cases:
  
\begin{parts}
    \item $ \text{GRR} = 2.10$, $ \text{NRR} = 1.10$
    \item $ \text{GRR} = 1.15$, $ \text{NRR} = 1.10$
    \item $ \text{GRR} = 1.75$, $ \text{NRR} = 1.80$
  \end{parts}\pts{7}
\end{parts}

\medskip\rule{\linewidth}{0.2pt}

\medskip


\noindent   \textbf{Question 2.}

\begin{parts}
  \item Discuss about the logistic curve used for measuring the growth of a population. \pts{7}
  \item What do you mean by a life table? Write the assumptions of a life table clearly. Discuss the various columns of a complete life table and also give the uses of a life table. \pts{7}
\end{parts}

\medskip\rule{\linewidth}{0.2pt}

\medskip


\noindent   \textbf{Question 3.}

\begin{parts}
  \item Define crude death rate and point out its limitations. Explain clearly the standardised death rates, including the methods of their computation. \pts{7}
  \item Describe the different measures of fertility with their merits and demerits. \pts{7}
\end{parts}

\medskip\rule{\linewidth}{0.2pt}

\medskip


\noindent   \textbf{Question 4.}

\begin{parts}
  \item Discuss the importance and limitations of index numbers. \pts{7}
  \item What is the need for a family budget enquiry in the construction of the cost of living index number? \pts{7}
\end{parts}

\medskip\rule{\linewidth}{0.2pt}

\medskip


\noindent   \textbf{Question 5.}

\begin{parts}
  \item Define the time reversal test. Show that Fisher's index number satisfies the time reversal test. \pts{7}
  \item What is statistical quality control (SQC)? Describe the main principles of SQC charts. Discuss the need and utility of SQC in industry. \pts{7}
\end{parts}

\medskip\rule{\linewidth}{0.2pt}

\medskip


\noindent   \textbf{Question 6.}

\begin{parts}
  \item Describe control charts for mean ($ar{X}$-chart) and standard deviation ($S$-chart). \pts{7}
  \item Discuss in detail the construction of $p$-chart and $c$-chart for attributes. \pts{7}
\end{parts}

\medskip\rule{\linewidth}{0.2pt}

\medskip


\noindent   \textbf{Question 7.}

\begin{parts}
  \item Define infant mortality rate. Why do we need to adjust the infant mortality rate? Describe in detail the Logan method of adjustment of the infant mortality rate. \pts{7}
  \item Explain secular trend in detail. Critically examine the various methods of measuring trend. \pts{7}
\end{parts}

\medskip\rule{\linewidth}{0.2pt}

\medskip


\noindent   \textbf{Question 8.}

\begin{parts}
  \item Write short notes on the different types of components of time series and explain the difference between seasonal and cyclical variations with examples. \pts{7}
  \item What do you mean by time series? Describe its uses in detail. \pts{7}
\end{parts}

\medskip\rule{\linewidth}{0.2pt}

\medskip


\noindent   \textbf{Question 9.}

\begin{parts}
  \item Explain the additive and multiplicative models of time series, stating clearly the assumptions underlying each model. \pts{7}
  \item Explain the moving average method of trend measurement. How does one choose the appropriate length of the moving average? \pts{7}
\end{parts}

\vfill
\rule{\linewidth}{0.4pt}

\begin{center}\small
  \href{https://bhu-exam.vercel.app/}{Click here to visit our website}\\[4pt]
  \href{https://me-aryan.vercel.app/}{Click here to visit our contributor}\\[4pt]
  \href{https://quashan-me.vercel.app/}{Click here to visit our contributor}
\end{center}

\end{document}